% =====================================================================
%  Formelsammlung: Crestfaktor bei N = 3 und N = 4 Toenen
%  (Gl. 109 ff. aus Crestfaktor_Randbedingung.pdf)
%  Braucht nur amsmath (align, \boxed, \tfrac, \operatorname).
%  Zum Einfuegen in ein bestehendes Dokument; keine Praeambel.
% =====================================================================

% ---------------------------------------------------------------------
\subsection*{A. Drei T\"one: Kr\"ummung und Crestfaktor}
% ---------------------------------------------------------------------

Kr\"ummung des Phasenprofils:
\begin{equation}
  \psi = \varphi_1 - \frac{\varphi_0+\varphi_2}{2}
  \label{eq:cf-psi}
\end{equation}

Amplituden und Einh\"ullende ($f_\mathrm{c}=\bar f=f_1$, Tonabstand $d$):
\begin{equation}
  a_0=a_2=\varrho:=\sqrt{r_x},\quad a_1=1,\qquad
  A(t) = \varrho\,e^{\mathrm{i}(\varphi_0-2\pi d t)} + e^{\mathrm{i}\varphi_1}
       + \varrho\,e^{\mathrm{i}(\varphi_2+2\pi d t)}
\end{equation}

Au\ss{}enphasen als Mittelwert plus Differenz:
\begin{equation}
  \varphi_0-\varphi_1 = -\psi+\delta,\qquad
  \varphi_2-\varphi_1 = -\psi-\delta,\qquad
  \delta = \tfrac12(\varphi_0-\varphi_2)
\end{equation}

\begin{align}
  A(t)\,e^{-\mathrm{i}\varphi_1}
  &= \varrho\,e^{\mathrm{i}(-\psi+\delta-2\pi d t)} + 1
   + \varrho\,e^{\mathrm{i}(-\psi-\delta+2\pi d t)}\\
  &= \varrho\,e^{-\mathrm{i}\psi}\Big[e^{\mathrm{i}(\delta-2\pi d t)}
   + e^{-\mathrm{i}(\delta-2\pi d t)}\Big] + 1\\
  &= 2\varrho\,e^{-\mathrm{i}\psi}\cos\Theta + 1,
  \qquad \Theta = 2\pi d\,(t-t_\delta),\quad t_\delta = \frac{\delta}{2\pi d}
\end{align}

\begin{equation}
  |A|^2 = 4\varrho^2\cos^2\Theta + 4\varrho\cos\psi\,\cos\Theta + 1
\end{equation}

Maximum bei $\cos\Theta=\operatorname{sign}(\cos\psi)$:
\begin{equation}
  \boxed{\;
  \max_t|A| = \sqrt{4\varrho^2+4\varrho\,|\cos\psi|+1},\qquad
  \mathrm{rms}|A| = \sqrt{2\varrho^2+1},\qquad
  C_x(\psi) = \sqrt2\,\frac{\sqrt{4\varrho^2+4\varrho\,|\cos\psi|+1}}{\sqrt{2\varrho^2+1}}
  \;}
  \label{eq:cf-crestpsi}
\end{equation}

\begin{align}
  \psi = 0 \text{ oder } \pi:&\quad \max|A| = 2\varrho+1 = \textstyle\sum_n a_n\\
  \psi = \pm\tfrac{\pi}{2}:&\quad \max|A| = \sqrt{4\varrho^2+1}
\end{align}

\begin{center}
\begin{tabular}{c|cccccc}
  \hline
  $\psi$ & $0^\circ$ & $45^\circ$ & $60^\circ$ & $90^\circ$ & $135^\circ$ & $180^\circ$\\
  \hline
  $C_x$ & $2{,}4494$ & $2{,}2836$ & $2{,}1586$ & $1{,}8220$ & $2{,}2836$ & $2{,}4494$\\
  \hline
\end{tabular}
\\[2pt]
{\small $\varrho=\sqrt{0{,}97}=0{,}98489$, $\mathrm{rms}|A|=1{,}71464$}
\end{center}

% ---------------------------------------------------------------------
\subsection*{B. Wovon das Maximum abh\"angt}
% ---------------------------------------------------------------------

Invarianzen von $\max_t|A|$:
\begin{align}
  \varphi_n &\to \varphi_n+\alpha: && A\to e^{\mathrm{i}\alpha}A,\\
  \varphi_n &\to \varphi_n+\delta n: && A(t)\to e^{\mathrm{i}\alpha'}A(t+\tau)
\end{align}

Zweite Differenzen (invariant unter beiden, $N-2$ St\"uck):
\begin{equation}
  \beta_k := \varphi_k - 2\varphi_{k+1} + \varphi_{k+2},\qquad k=0,\dots,N-3
  \label{eq:cf-beta}
\end{equation}
\begin{equation}
  \max_t|A| = F(\beta_0,\dots,\beta_{N-3}),\qquad
  N=3:\ \beta_0 = -2\psi
\end{equation}

Mit Zentrierung $\varphi_n+\varphi_{N-1-n}=c$:
\begin{align}
  N=3:&\quad \beta_0 = (\varphi_0+\varphi_2)-2\varphi_1 = c-2\varphi_1 \in\{0,\,2\pi\}\equiv 0\\
  N=4:&\quad \beta_0 = \varphi_0-2\varphi_1+\varphi_2,\qquad
         \beta_1 = \varphi_1-2\varphi_2+\varphi_3 = -\beta_0
\end{align}
\begin{equation}
  \varphi_y = [\,0,\ b,\ c-b,\ c\,]:\qquad \beta_0 = c-3b,\qquad \beta_1 = 3b-c
\end{equation}
\begin{equation}
  \text{Arbeitspunkt } c=114^\circ,\ b=98^\circ:\qquad
  \beta_0 = -180^\circ,\quad \beta_1 = +180^\circ
\end{equation}

% ---------------------------------------------------------------------
\subsection*{C. Vier T\"one: die Einh\"ullende der zentrierten Familie}
% ---------------------------------------------------------------------

\noindent\textit{a) T\"one relativ zu $f_\mathrm{c}=\bar f$:}
\begin{equation}
  f_n = f_\mathrm{off}+n f_\mathrm{env},\quad
  \bar f = f_\mathrm{off}+\tfrac32 f_\mathrm{env},\qquad
  f_n-f_\mathrm{c} = \bigl(n-\tfrac32\bigr)f_\mathrm{env} =: \hat n\,f_\mathrm{env},\quad
  \hat n\in\bigl\{-\tfrac32,-\tfrac12,+\tfrac12,+\tfrac32\bigr\}
\end{equation}

\noindent\textit{b) Zeitvariable:}
\begin{equation}
  u := \pi f_\mathrm{env} t
  \quad\Longrightarrow\quad
  2\pi\hat n f_\mathrm{env} t = 2\hat n u \in\{-3u,\,-u,\,+u,\,+3u\},
  \qquad t\in[0,T_\mathrm{env})\ \leftrightarrow\ u\in[0,\pi)
\end{equation}

\noindent\textit{c) Phasen zerlegen:}
\begin{equation}
  \chi_n := \varphi_n-\tfrac{c}{2}
  \quad\Longrightarrow\quad
  \chi_n+\chi_{3-n} = 0,\qquad
  \chi_\mathrm{a} := \chi_3 = -\chi_0,\qquad
  \chi_\mathrm{i} := \chi_2 = -\chi_1
\end{equation}

\begin{center}
\begin{tabular}{c|c|c|c|c|c}
  \hline
  $n$ & $\hat n$ & $a_n$ & $\varphi_n$ & $\chi_n$ & Arbeitspunkt\\
  \hline
  0 & $-\tfrac32$ & $\varrho$ & $0$   & $-c/2 = -\chi_\mathrm{a}$  & $-57^\circ$\\
  1 & $-\tfrac12$ & $1$       & $b$   & $b-c/2 = -\chi_\mathrm{i}$ & $+41^\circ$\\
  2 & $+\tfrac12$ & $1$       & $c-b$ & $c/2-b = \chi_\mathrm{i}$  & $-41^\circ$\\
  3 & $+\tfrac32$ & $\varrho$ & $c$   & $c/2 = \chi_\mathrm{a}$    & $+57^\circ$\\
  \hline
\end{tabular}
\end{center}

\noindent\textit{d) Einsetzen, $e^{\mathrm{i}c/2}$ ausklammern:}
\begin{equation}
  a_n e^{\mathrm{i}(2\hat n u+\varphi_n)} = e^{\mathrm{i}c/2}\,a_n e^{\mathrm{i}(2\hat n u+\chi_n)}
\end{equation}
\begin{equation}
  A(t)\,e^{-\mathrm{i}c/2}
  = \varrho\,e^{\mathrm{i}(-3u-\chi_\mathrm{a})}
  + e^{\mathrm{i}(-u-\chi_\mathrm{i})}
  + e^{\mathrm{i}(+u+\chi_\mathrm{i})}
  + \varrho\,e^{\mathrm{i}(+3u+\chi_\mathrm{a})}
\end{equation}

\noindent\textit{e) Paarweise zusammenfassen ($e^{\mathrm{i}x}+e^{-\mathrm{i}x}=2\cos x$):}
\begin{align}
  \varrho\bigl[e^{-\mathrm{i}(3u+\chi_\mathrm{a})}+e^{+\mathrm{i}(3u+\chi_\mathrm{a})}\bigr]
  &= 2\varrho\cos(3u+\chi_\mathrm{a})\\
  e^{-\mathrm{i}(u+\chi_\mathrm{i})}+e^{+\mathrm{i}(u+\chi_\mathrm{i})}
  &= 2\cos(u+\chi_\mathrm{i})
\end{align}
\begin{equation}
  A(t)\,e^{-\mathrm{i}c/2} = 2\varrho\cos(3u+\chi_\mathrm{a}) + 2\cos(u+\chi_\mathrm{i})
  \quad\text{(reell)}
\end{equation}
\begin{equation}
  \text{Arbeitspunkt:}\qquad |A| = \bigl|\,2\varrho\cos(3u+57^\circ)+2\cos(u-41^\circ)\,\bigr|
\end{equation}

\noindent\textit{f) Zeitnullpunkt verschieben, $u' := u+\chi_\mathrm{i}$:}
\begin{align}
  u+\chi_\mathrm{i} &= u',\\
  3u+\chi_\mathrm{a} &= 3u'-3\chi_\mathrm{i}+\chi_\mathrm{a} = 3u'+\gamma
\end{align}
\begin{equation}
  \boxed{\;
  A\,e^{-\mathrm{i}c/2} = 2\varrho\cos(3u+\gamma)+2\cos u,\qquad
  \gamma = \chi_\mathrm{a}-3\chi_\mathrm{i}
         = \tfrac{c}{2}-3\bigl(\tfrac{c}{2}-b\bigr) = 3b-c = -\beta_0
  \;}
  \label{eq:cf-n4env}
\end{equation}
\begin{equation}
  \gamma = 57^\circ-3\cdot(-41^\circ) = 3\cdot98^\circ-114^\circ = 180^\circ
\end{equation}
\begin{equation}
  u\to u+\pi:\quad \cos u\to-\cos u,\quad \cos(3u+\gamma)\to-\cos(3u+\gamma)
  \quad\Longrightarrow\quad |A|\ \text{hat Periode}\ \pi
\end{equation}

\noindent\textit{g) Invarianz von $\gamma$ unter einer Rampe $\varphi_n\to\varphi_n+\delta n$:}
\begin{align}
  &c\to c+3\delta,\qquad \chi_n\to\chi_n+\delta\hat n,\qquad
  \chi_\mathrm{a}\to\chi_\mathrm{a}+\tfrac32\delta,\quad
  \chi_\mathrm{i}\to\chi_\mathrm{i}+\tfrac12\delta,\\
  &\gamma\to\gamma+\tfrac32\delta-3\cdot\tfrac12\delta = \gamma
\end{align}

% ---------------------------------------------------------------------
\subsection*{D. Vier T\"one: Grenzf\"alle und geschlossene Form}
% ---------------------------------------------------------------------

Rampe, $\gamma=0$:
\begin{align}
  A\,e^{-\mathrm{i}c/2} &= 2\varrho\cos3u+2\cos u,\\
  \max|A| &= 2\varrho+2 = \textstyle\sum_n a_n = 4{,}1540659,\qquad
  C_y = \sqrt2\,\frac{4{,}1540659}{2{,}0784610} = 2{,}8264839
\end{align}

Arbeitspunkt, $\gamma=\pi$:
\begin{equation}
  A\,e^{-\mathrm{i}c/2} = 2\cos u - 2\varrho\cos3u,\qquad
  u=0:\ 2-2\varrho
\end{equation}

\noindent\textit{h) Dreifachwinkel und Polynom:}
\begin{equation}
  \cos3u = \cos2u\cos u-\sin2u\sin u
         = (2\cos^2u-1)\cos u-2\sin^2u\cos u
         = 4\cos^3u-3\cos u
\end{equation}
\begin{equation}
  v:=\cos u\in[-1,1]:\qquad
  G(v) := A\,e^{-\mathrm{i}c/2} = 2v-2\varrho(4v^3-3v) = (2+6\varrho)\,v-8\varrho\,v^3
\end{equation}

\noindent\textit{i) Symmetrie und Kandidaten:}
\begin{equation}
  G(-v)=-G(v),\qquad
  G(0)=0,\qquad
  G(1)=2+6\varrho-8\varrho=2-2\varrho,\qquad
  G'(v)=(2+6\varrho)-24\varrho\,v^2
\end{equation}

\noindent\textit{j) Extremstelle:}
\begin{equation}
  G'(v_*)=0:\quad v_*^2 = \frac{2+6\varrho}{24\varrho} = \frac{1+3\varrho}{12\varrho},
  \qquad G''(v)=-48\varrho v<0,
  \qquad v_*\le1\ \Leftrightarrow\ \varrho\ge\tfrac19
\end{equation}

\noindent\textit{k) Wert:}
\begin{equation}
  G(v_*) = v_*\bigl[(2+6\varrho)-8\varrho v_*^2\bigr],\qquad
  8\varrho v_*^2 = 8\varrho\cdot\frac{1+3\varrho}{12\varrho} = \tfrac23(1+3\varrho)
\end{equation}
\begin{equation}
  (2+6\varrho)-\tfrac23(1+3\varrho) = 2(1+3\varrho)-\tfrac23(1+3\varrho) = \tfrac43(1+3\varrho)
\end{equation}
\begin{equation}
  G(v_*) = \tfrac43(1+3\varrho)\,v_*
         = \tfrac43(1+3\varrho)\sqrt{\frac{1+3\varrho}{12\varrho}}
         = \tfrac43\,\frac{(1+3\varrho)^{3/2}}{2\sqrt{3\varrho}}
         = \tfrac23\,\frac{(1+3\varrho)^{3/2}}{\sqrt{3\varrho}}
\end{equation}
\begin{equation}
  \boxed{\;
  \max_t|A| = \frac23\,\frac{(1+3\varrho)^{3/2}}{\sqrt{3\varrho}},\qquad
  C_y = \sqrt2\;\frac{\tfrac23(1+3\varrho)^{3/2}/\sqrt{3\varrho}}{\sqrt{2\varrho^2+2}}
  \;}
  \label{eq:cf-n4closed}
\end{equation}

\noindent\textit{l) Zahlen, $\varrho=\sqrt{1{,}16}=1{,}0770330$:}
\begin{center}
\begin{tabular}{l|r}
  \hline
  $1+3\varrho$ & $4{,}2310991$\\
  $12\varrho$ & $12{,}9243955$\\
  $v_*^2$ & $0{,}3273731$\\
  $v_*=\cos u_*$ & $0{,}5721652$\quad ($u_*=55{,}10^\circ$, verschobene Zeit)\\
  $\tfrac43(1+3\varrho)$ & $5{,}6414655$\\
  $\max|A| = 5{,}6414655\cdot0{,}5721652$ & $3{,}2278503$\\
  $|G(1)| = |2-2\varrho|$ & $0{,}1540659$\\
  $\sqrt{\sum a_n^2}=\sqrt{2r_y+2}=\sqrt{4{,}32}$ & $2{,}0784610$\\
  $C_y = \sqrt2\cdot3{,}2278503/2{,}0784610$ & $2{,}1962739$\\
  \hline
\end{tabular}
\end{center}

\noindent\textit{m) Zur\"uck in die echte Zeit ($u=u'-\chi_\mathrm{i}=u'+41^\circ$, $t=(u/180^\circ)\,T_\mathrm{env}$):}
\begin{align}
  v=+v_*:&\quad u'=55{,}10^\circ,\quad u=96{,}10^\circ,\quad t=4{,}3288\ \mu\mathrm{s},\\
  v=-v_*:&\quad u'=124{,}90^\circ,\quad u=165{,}90^\circ,\quad t=7{,}4730\ \mu\mathrm{s}
\end{align}
